\section{The Value of Profiling Evidence}
\label{sec:evidence}
% NEW (written for the TOMPECS version; not in the original paper)
This section addresses RQ4 in two steps: whether profiling data helps at all (an arm with and an
arm without it), and what happens when the data is converted into named remedies before the model
sees it.

\subsection{Does Profiling Data Help?}
\label{sec:profabl}
Section~\ref{sec:ablation} ablates components \emph{within} the prompt while always supplying profiling
data. The prior question---whether the profiling data helps at all---requires an arm with no profile
data, which we add here.

\paragraph{Arms.}
The \texttt{pcrl} arm is the standard prompt carrying PC-sampling and roofline summaries only (no
derived kernel-budget, occupancy or diagnostic summaries). The \texttt{code\_only} arm removes the
entire performance-analysis block. Removing it alone would bias the comparison, because four
instruction sentences refer to the data---most importantly the rule \textit{``if no safe substantive
optimization can be justified from the provided performance data, keep the code semantically
unchanged''}, which would invite a no-op when no data is present. Those four sentences are therefore
reworded to refer to the model's analysis of the source; the two prompts are otherwise byte-identical
(verified: the diff is the removed block plus those four sentences). Both arms were regenerated with
the same code base, and all pairs use the same correctness gates and interleaved measurement as
Section~\ref{sec:realapps}.

\paragraph{Scoring.}
HeCBench pairs are measured on a shared H100 node where per-repetition spread is sometimes larger than
any plausible effect, so we score a pair as a \emph{tie} when $|\mathrm{pcrl}/\mathrm{co}-1|$ is smaller
than the largest repetition spread measured for that pair---that is, when the difference between the
arms does not exceed its own measurement noise. Two applications (\texttt{boxfilter},
\texttt{ccsd-trpdrv}) are excluded because the \emph{unmodified} program exceeds our $600$\,s per-run
cap on that node, and one (\texttt{background-subtract}, GPT-5) because the \texttt{code\_only} arm
produced no correct candidate at all, which we report under correctness rather than speed. One
caveat carries through the whole benchmark half of this comparison: on this node the benchmarks' own
verification gives no usable verdict for the \emph{unmodified} program---five of nine report FAIL and
the rest report nothing---so a candidate's verdict is informative only where it differs from the
pristine one (Section~\ref{sec:limits}).

\paragraph{Result.}
Profiling data helps on the large real application and makes no measurable difference on the small
benchmarks (Table~\ref{tab:profabl}). On Cholla both models are clearly faster with the profile
($+11.5\%$ for GPT-5, $+7.7\%$ for Gemini), and a second, independently generated \texttt{pcrl} sample
keeps both on the same side of the comparison (GPT-5 $1.187$, Gemini $1.023$ relative to
\texttt{code\_only}), although the Gemini replicate is much weaker than its first sample. On
AutoDock-GPU both arms reach the same ${\sim}1.03\times$. All seventeen scored HeCBench pairs are ties
by the rule above: the largest arm-to-arm difference is $7.8\%$ for GPT-5 (\texttt{adv}, inside a
$31.8\%$ spread) and $33.2\%$ for Gemini (\texttt{background-subtract}, inside a $147.6\%$ spread), and
the eight pairs whose spread is under $5\%$ all land between $0.998$ and $1.003$. Not one benchmark
shows a difference that its own measurement noise does not already cover.


\begin{table}[t]
\centering
\footnotesize
\caption{Profile-data ablation: speedup against the unmodified application for the prompt without
  profile data (\texttt{co}) and with PC sampling + roofline (\texttt{pcrl}), and their ratio.
  A HeCBench pair is a tie when the arm-to-arm difference is smaller than that pair's repetition
  spread; the entry gives the largest difference observed and the spread it sits inside.
  Application medians: Cholla over $4$ of $7$ repetitions (three excluded for drift with no foreign
  job detected), AutoDock-GPU over $2$ of $7$ (five excluded for foreign jobs).}
\label{tab:profabl}
\setlength{\tabcolsep}{2pt}
\begin{tabular}{@{}llrrr@{}}
\toprule
\textbf{Application} & \textbf{Model} & \textbf{co} & \textbf{pcrl} & \textbf{pcrl/co} \\
\midrule
\multirow{2}{*}{Cholla}
  & GPT-5  & 1.324$\times$ & \textbf{1.476}$\times$ & \textbf{1.115} \\
  & Gemini & 1.471$\times$ & \textbf{1.584}$\times$ & \textbf{1.077} \\
\midrule
\multirow{2}{*}{AutoDock-GPU}
  & GPT-5  & 1.030$\times$ & 1.031$\times$ & 1.001 \\
  & Gemini & 1.035$\times$ & 1.035$\times$ & 1.000 \\
\midrule
HeCBench (8 pairs) & GPT-5  & \multicolumn{3}{c}{all ties} \\
HeCBench (9 pairs) & Gemini & \multicolumn{3}{c}{all ties} \\
\bottomrule
\end{tabular}
\end{table}

\paragraph{Correctness, not only speed.}
Two HeCBench pairs favour the profile arm on correctness rather than time. On
\texttt{background-subtract}, GPT-5 without profile data failed the application's own output check in
all three repair iterations and produced no usable candidate, while the \texttt{pcrl} arm produced a
correct one; on \texttt{atomicSystemWide}, the \texttt{code\_only} candidate fails verification and the
\texttt{pcrl} candidate does not.

\paragraph{Interpretation.}
The pattern is consistent with the workloads rather than with the models. A HeCBench kernel is a few
hundred lines with one obvious hot loop, and the profile tells a model little it cannot infer from the
source. Cholla's HLLC solver is 307 lines inside a 21k-line CUDA application whose cost is spread over
several kernels; there the profile identifies both the kernel worth editing and the reason it is slow
(latency-bound at $19.7\%$ DRAM and $40.4\%$ SM throughput, occupancy $23\%$ limited by 109 registers
per thread). We therefore expect profile data to pay off in proportion to how hard the bottleneck is to
locate from source alone---which is precisely the regime real applications occupy.
A separate question is what happens when that data is not merely supplied but \emph{interpreted} into
named remedies before the model sees it; Section~\ref{sec:guidvar} measures that arm.

\subsection{When Profiling Data Becomes Prescriptive}
\label{sec:guidvar}
Sections~\ref{sec:ablation} and \ref{sec:profabl} treat profiling data as something the prompt either
carries or does not. A third option is to carry it \emph{interpreted}: to convert the metrics into
named remedies before the model sees them. Recent work argues strongly for this---KernelPro turns
Nsight Compute, SASS and Nsight Systems output into natural-language guidance through a set of
micro-profiling tools, and reports a $125\%$ higher speedup than handing the model raw metrics
\cite{gai2026kernelpro}. We evaluated the same idea, as a \texttt{diagnostics} prompt that adds a
bounded checklist of measured conditions and the transformation each one suggests, and reach a more
qualified conclusion: interpretation does not shift the average outcome so much as it widens the
distribution of outcomes.

\paragraph{Setup.}
% EDIT (2026-10-26): chemv removed (see the last paragraph of this section); concat/Gemini and
% boxfilter/GPT-5 effects re-attributed to the structural prompt; the >5x boxfilter results inspected.
Table~\ref{tab:guidvar} reports four arms on one node (GH200, idle, all arms built once and timed
round-robin in a single session): the unmodified program, the constrained \emph{baseline} prompt, the
relaxed \emph{structural} prompt, and \emph{diagnostics}, which is the structural prompt plus the
checklist. The cleanest comparison is therefore \emph{diag} against \emph{struct}, which differ only
in the checklist. Half of the rows have a repetition spread below $5\%$, including the two largest
effects of the checklist, so those differences are not measurement artifacts; each cell is, however,
a single generation (Section~\ref{sec:limits}).

\begin{table}[t]
\centering
\footnotesize
\setlength{\tabcolsep}{2pt}
\caption{Effect of converting profiling metrics into a prescriptive checklist (\emph{diag}), against
  the constrained (\emph{base}) and relaxed (\emph{struct}) prompts. Speedups are interleaved medians
  against the unmodified program on one GH200. \textbf{Bold} marks the cases discussed in the text.
  ``spread'' is the largest repetition spread across the four arms. \texttt{chemv} is excluded
  (see text).}
\label{tab:guidvar}
\resizebox{\columnwidth}{!}{%
\begin{tabular}{@{}llrrrr@{}}
\toprule
\textbf{Benchmark} & \textbf{Model} & \textbf{base} & \textbf{struct} & \textbf{diag} & \textbf{spread} \\
\midrule
\texttt{concat}               & GPT-5  & 1.031$\times$ & 0.838$\times$ & \textbf{0.349}$\times$ & 0.1\% \\
\texttt{boxfilter}            & GPT-5  & \textbf{10.280}$\times$ & 6.352$\times$ & 5.503$\times$ & 0.0\% \\
\midrule
\texttt{bincount}             & GPT-5  & 1.035$\times$ & 1.037$\times$ & \textbf{1.374}$\times$ & 5.2\% \\
\texttt{bincount}             & Gemini & 1.106$\times$ & 1.003$\times$ & \textbf{1.380}$\times$ & 4.4\% \\
\texttt{concat}               & Gemini & 1.001$\times$ & \textbf{1.797}$\times$ & 1.798$\times$ & 46.8\% \\
\texttt{boxfilter}            & Gemini & 8.691$\times$ & 8.930$\times$ & 9.371$\times$ & 42.6\% \\
\texttt{convolutionSeparable} & Gemini & 1.418$\times$ & 1.499$\times$ & 1.514$\times$ & 19.3\% \\
\midrule
\texttt{blockAccess}          & both   & $\approx 1.00\times$ & $\approx 1.00\times$ & $\approx 1.00\times$ & 0.2\% \\
\bottomrule
\end{tabular}%
}
\end{table}

\paragraph{The checklist produces our largest gain.}
On \texttt{bincount} it takes both models from roughly $1.00$--$1.04\times$ under the structural prompt
to $1.37$--$1.38\times$, approaching our hand-tuned reference ($1.392\times$). Where the measured
condition is real and the named remedy is the right one, telling the model what to do works. The
other large improvement in the table is not the checklist's: on \texttt{concat} under Gemini the
$1.80\times$ appears already under the structural prompt ($1.797\times$), whose relaxed syntactic
rules permit the vectorized templating it needs, and the checklist adds nothing ($1.798\times$) at a
$46.8\%$ spread.

\paragraph{It also produces our largest loss.}
On \texttt{concat}, GPT-5 followed the instruction to process several elements per thread literally,
to one-thread-per-128-element line: a $128\times$ collapse in parallelism, measured at $0.349\times$
against $0.838\times$ without the checklist, with a $0.1\%$ spread. On \texttt{boxfilter}, GPT-5's
$10.280\times$ under the baseline prompt falls to $6.352\times$ when the syntactic rules are relaxed and
to $5.503\times$ with the checklist added; most of that loss belongs to the relaxation, not to the
checklist. All \texttt{boxfilter} results exceed our $5\times$ alarm, so we inspected the diffs: every
candidate keeps the timing loop and both kernel launches per cycle, and gains by tiling the column
pass over rows, which the original runs as one thread per image column ($30$ blocks of $64$ threads
for a $1920$-pixel-wide image); the benchmark's own check compares every output pixel exactly with a
host implementation.

\paragraph{Interpretation narrows the search.}
A checklist of remedies is an anchor as much as a hint. When the listed transformation is the best
available one, anchoring is worth a large factor; when it is not, the model applies it anyway, as
GPT-5 did on \texttt{concat}. Averaged over all twelve (benchmark, model) pairs we measured the effect
is close to neutral---a geometric mean of $0.97\times$ against the structural prompt, with
individual pairs from $0.42\times$ to $1.38\times$---which is why a mean speedup, the statistic the
literature reports, hides it. The dispersion is the result: the same prompt component that produces
$1.380\times$ on one benchmark produces $0.349\times$ on another, both at sub-$5\%$ spread.

Two practical consequences follow. First, prescriptive guidance should be reported with its variance,
not only its mean; we would not have seen this in a geometric-mean summary. Second, the guidance should
be \emph{evidence rather than remedy}---state the measured condition (``$23\%$ occupancy, limited by
109 registers per thread'') and let the model decide what follows, instead of naming the
transformation. We did not get to test that variant, and it is the first thing we would run next.

\paragraph{An excluded outlier: \texttt{chemv}.}
% NEW (2026-10-26). Evidence: chemv-verification-vacuous (experiments memory), report sec. 13.3/14,
% and a re-check on illyad (CUDA 13.3.1, 2026-10-26) of every chemv candidate with REPEAT=1 against
% the CPU reference: max relative error 2.2e-7 for the unmodified program and for Gemini's baseline
% candidate, at most 1.1e-6 for any candidate.
We exclude \texttt{chemv} from every table because its built-in verification cannot certify a result.
Its timing loop launches both kernels $1000$ times, accumulating into the same output vector without
re-initializing it, and then compares that vector with a single application of the CPU reference: the
check fails for the unmodified program, or---once the accumulated values overflow to NaN---passes
vacuously, because a comparison with NaN is never true. The case is nonetheless instructive. Under the
baseline prompt, Gemini rewrote both polyhedral-compiler-generated kernels: it replaced each thread's
tiled loop nest (an outer loop over $32$-element tiles with a barrier per tile, and an inner loop
clipped to the tile) by a single loop over the same elements of its row, removed the barriers---which
the computation does not need, since each thread updates only its own output element and the kernels
use no shared memory---and accumulated in registers instead of reading and writing the output in global
memory on every iteration. The number of multiply-adds is unchanged; what was removed is loop structure
that correct output does not require. The benchmark timed it at $7.8\times$, over our $5\times$ alarm,
so we checked it outside the benchmark's harness, running one application of the two kernels (no
repeat loop, so nothing overflows) against the CPU reference: its largest relative error is
$2.2\times10^{-7}$, identical to the unmodified program's, and no \texttt{chemv} candidate from either
model in the baseline, structural, checklist or profile-data arms exceeds $1.1\times10^{-6}$. Under the checklist the same model kept the original tile
loops and added the shared-memory staging the checklist named. We report the case as an example of the
pipeline removing loops that correctness does not need---the legitimate counterpart of the gaming of
Section~\ref{sec:intent}, which removed a loop the measurement does need---and not as a measured
speedup, since the timing that produced it rests on a loop whose output is NaN.

